\section{Recent-path Decomposition}

This section translates residual signals into price-path structures. Its first three subsections use the three-model CNN--MLP-4--HistGB decomposition defined in Section 6; hence unqualified CNN-specific and MLP-4-specific refer only to that estimand. The final subsection separately analyzes the pairwise temporal CNN--1D-CNN residuals and labels each component with its conditioning model.

\subsection{Lag Sensitivity}

Intervening on individual days in the most recent five-day window reveals sharp differences in the three-model components. The CNN-specific signal is highly sensitive to particular near-end lags, especially $t-1$, $t$, and $t-3$. In contrast, removing individual tabular days does not reduce the MLP-4-specific Sharpe, indicating that its signal is not concentrated in a single recent observation.

\begin{table}[htbp]
\centering
\caption{Recent-lag sensitivity of specific signals}
\label{tab:recent_lag}
\begin{tabular}{lrrrr}
\toprule
Lag & CNN $\Delta$Sharpe & CNN drop & MLP $\Delta$Sharpe & MLP drop \\
\midrule
$t-4$ & 0.07 & 5\% & -0.29 & -14\% \\
$t-3$ & 0.61 & 48\% & -0.46 & -22\% \\
$t-2$ & -0.15 & -12\% & -0.20 & -9\% \\
$t-1$ & 0.79 & 62\% & -0.04 & -2\% \\
$t$ & 0.62 & 49\% & -0.03 & -1\% \\
\bottomrule
\end{tabular}
\begin{minipage}{0.92\linewidth}
\footnotesize
\textit{Notes:} Positive $\Delta$Sharpe means the intervention reduces the specific signal's portfolio Sharpe. CNN interventions flatten the corresponding image day; MLP interventions remove the corresponding tabular day and retrain.
\end{minipage}
\end{table}


\subsection{Path Shape Beyond Recent Return}

Controlling for cumulative five-day return does not eliminate path-shape explanatory power. Variables such as maximum drawdown, path volatility, late-versus-early return, reversal, and monotonicity add substantial explanatory power for specific scores. The three-model CNN-specific component loads more on drawdown and path-volatility structures, while the MLP-4-specific component is more closely aligned with monotonicity and continuation-like geometry.

\begin{table}[htbp]
\centering
\caption{Path-shape explanatory power beyond five-day return}
\label{tab:recent_pathshape}
\begin{tabular}{lrrrr}
\toprule
Score & $R^2$(CumRet5) & $R^2$(CumRet5 + Path) & $\Delta R^2$ & NW $t$ \\
\midrule
CNN-specific & 0.021 & 0.072 & 0.052 & 45.6 \\
MLP-specific & 0.048 & 0.089 & 0.041 & 36.6 \\
\bottomrule
\end{tabular}
\begin{minipage}{0.92\linewidth}
\footnotesize
\textit{Notes:} Path variables include late-minus-early return, path volatility, up-day count, maximum drawdown, reversal, and monotonicity over the most recent five days.
\end{minipage}
\end{table}


\subsection{Component Sensitivity}

Matched interventions show that the CNN-specific signal is most sensitive to volume removal. MLP-4-specific performance is not reduced by any single engineered component block, reinforcing the interpretation that its signal is distributed across tabular geometry rather than tied to one component.

\begin{table}[htbp]
\centering
\caption{Matched component sensitivity}
\label{tab:recent_blocks}
\begin{tabular}{lrrrr}
\toprule
Component & CNN $\Delta$Sharpe & CNN drop & MLP $\Delta$Sharpe & MLP drop \\
\midrule
Volume & 0.36 & 28\% & -0.26 & -12\% \\
Wick & 0.07 & 5\% & -0.10 & -4\% \\
Direction & 0.03 & 2\% & 0.09 & 4\% \\
Return & 0.02 & 1\% & -0.28 & -13\% \\
Gap & -0.25 & -19\% & -0.45 & -21\% \\
Intraday & 0.07 & 5\% & -0.23 & -11\% \\
\bottomrule
\end{tabular}
\begin{minipage}{0.92\linewidth}
\footnotesize
\textit{Notes:} Positive $\Delta$Sharpe means that removing the component reduces the specific-score Sharpe. Negative values indicate an improvement after removal.
\end{minipage}
\end{table}


Overall, the specific signals are not reducible to a simple recent-return effect. They represent different ways of parsing the same recent price path.

\subsection{CNN versus Temporal CNN}

The temporal benchmark makes the mechanism comparison more stringent because both models impose local convolutional structure on the same 20-day history. Table~\ref{tab:temporal_recent_path} reports the parallel E4T analysis using the \emph{pairwise temporal} residuals, not the three-model components above. For both pairwise residuals, adding recent-path shape variables beyond CumRet5 materially increases daily cross-sectional explanatory power. The CNN-specific component conditional on 1D-CNN is especially sensitive to $t-1$, $t-3$, and volume removal. The 1D-CNN-specific component conditional on CNN is instead most sensitive to removal of the recent return path, which reduces its Sharpe by 1.01. These are intervention-based model-sensitivity results, not causal effects of individual chart variables, but they show that residual disagreement is not merely a difference in the strength of one temporal signal.

\begin{table}[htbp]
\centering
\caption{Recent-path mechanisms: pairwise CNN--1D-CNN specific signals}
\label{tab:temporal_recent_path}
\resizebox{\textwidth}{!}{%
\begin{tabular}{lrrrl}
\toprule
Component & $R^2$(CumRet5) & $R^2$(CumRet5 + shape) & $\Delta R^2$ & Main intervention sensitivity \\
\midrule
CNN-specific conditional on 1D-CNN & 0.0159 & 0.0792 & 0.0633 & $t-1$: 0.711; $t-3$: 0.600; volume: 0.202 \\
1D-CNN-specific conditional on CNN & 0.0673 & 0.1352 & 0.0680 & return path: 1.006; $t-1$: 0.285 \\
\bottomrule
\end{tabular}%
}
\vspace{0.25em}
\footnotesize\textit{Notes:} Shape variables are LateEarly, PathVol, UpCount5, MaxDD5, Reversal5, and Monotonicity. Intervention entries are reductions in EW H--L Sharpe. Each $\Delta R^2$ has Newey--West $t>43$.
\end{table}

